\chapter*{Appendix S: Geometric Semantics of Modal Operators}
\addcontentsline{toc}{chapter}{Appendix S: Geometric Semantics of Modal Operators}

\section{Overview}

This appendix develops the geometric semantics underlying the modal operators introduced previously. Whereas the syntactic and Kripke-style logics formalised earlier describe invariants of reasoning, the present appendix interprets these modalities in the RSVP semantic manifold and in the Polyxan hypergraph fibration.  

The goal is to provide a geometric realisation of necessity predicates: a formula is necessary in a world precisely when it is stable under the canonical geometric transformations tied to curvature descent and rewrite confluence.

\section{RSVP Curvature Frames}

\subsection{Semantic Worlds}

Each world is interpreted as a smoothly embedded field configuration:
\[
w = (\Phi,\mathbf{v},S,\iota)
\]
where $\iota : \mathrm{Inc}(\mathcal{G}) \to \mathcal{M}$ is curvature-matched and harmonic.

\subsection{Accessibility}

The accessibility relation governing the RSVP modality is defined geometrically:
\[
w \;R_{\mathrm{RSVP}}\; w'
\quad\text{iff}\quad
\mathcal{E}_{\mathrm{total}}[w'] \leq \mathcal{E}_{\mathrm{total}}[w]
\quad\text{and}\quad
w' \text{ is reachable by a single curvature descent step.}
\]

Thus RSVP necessity expresses invariance under all admissible energy-decreasing flows.

\subsection{Curvature Persistence}

For any world $w$ and formula $A$,
\[
w \Vdash \Box_{\mathrm{RSVP}} A
\quad\Longleftrightarrow\quad
A \text{ holds in every curvature-descendant of } w.
\]

This matches the monotonicity of the energy functional and the stability of harmonic strata.

\section{Polyxan Rewriting Fibrations}

\subsection{Fibre Worlds}

A Polyxan world is an EHR-normal hypergraph:
\[
w = \mathcal{G}_{\mathrm{nf}}.
\]

\subsection{Rewriting Relation}

The rewriting relation is:
\[
\mathcal{G} \;R_{\mathrm{Polyx}}\; \mathcal{G}'
\quad\text{iff}\quad
\mathcal{G}' 
\text{ is obtained by an atomic EHR rewrite }
\text{and both lie in the same normalisation fibre.}
\]

This ensures access is limited to confluence-preserving transformations.

\subsection{Rewrite Necessity}

A formula holds necessarily when it is invariant under the entire confluence hull:
\[
\mathcal{G} \Vdash \Box_{\mathrm{Polyx}} A
\quad\Longleftrightarrow\quad
A \text{ holds in every } \mathcal{G}' \;\text{with}\; \mathcal{G}R_{\mathrm{Polyx}}^\ast \mathcal{G}'.
\]

\section{Joint Geometric Models}

\subsection{Product Worlds}

Semantic worlds are pairs:
\[
(w_{\mathrm{RSVP}},w_{\mathrm{Polyx}}).
\]

\subsection{Mixed Accessibility}

The mixed accessibility relation is:
\[
(w_1,w_2) 
\;R\;
(w_1',w_2')
\quad\Longleftrightarrow\quad
w_1 R_{\mathrm{RSVP}} w_1'
\quad\text{and}\quad
w_2 R_{\mathrm{Polyx}} w_2'.
\]

This yields a product frame for joint modalities.

\section{Categorical Semantics}

\subsection{World Category}

Define $\mathbf{W}$ as the category whose:
\begin{itemize}
\item objects are worlds,
\item morphisms are energy- or rewrite-nonincreasing transformations,
\item composition is concatenation of descent or rewrite steps.
\end{itemize}

\subsection{Modal Functors}

The two modalities are interpreted as endofunctors:
\[
\Box_{\mathrm{RSVP}}, \Box_{\mathrm{Polyx}} : \mathbf{W} \to \mathbf{W}.
\]

Each maps a world to the full subcategory of its descendants or rewrite variants.

\subsection{Geometric Natural Transformations}

There are natural transformations:
\[
\Box_{\mathrm{RSVP}} \Longrightarrow \mathrm{Id},
\qquad
\Box_{\mathrm{Polyx}} \Longrightarrow \mathrm{Id},
\]
corresponding to the reflexivity axioms.

\subsection{Commutativity}

Curvature descent and rewrite confluence commute up to natural isomorphism:
\[
\Box_{\mathrm{RSVP}}\Box_{\mathrm{Polyx}}
\;\cong\;
\Box_{\mathrm{Polyx}}\Box_{\mathrm{RSVP}}.
\]

\section{Sheaf Models}

\subsection{Base Site}

Let $\mathcal{B}$ be the site whose objects are semantic open sets in $\mathcal{M}$ and rewrite-stable patches of hypergraphs.

\subsection{Presheaves of Truth Values}

Truth values are interpreted as presheaves:
\[
\mathcal{F} : \mathcal{B}^{\mathrm{op}} \to \mathbf{Set},
\]
where $\mathcal{F}(U)$ is the set of formulas true in region $U$.

\subsection{Sheaf Condition}

The sheaf condition ensures global truth is determined by consistent local truth across overlaps, reflecting the stratified nature of the semantic manifold.

\subsection{Modalities as Sheaf Operators}

Each modality acts on presheaves via:
\[
(\Box_{\mathrm{RSVP}}\mathcal{F})(U)
  = \bigcap_{U' \subseteq U\,:\,U'\ R_{\mathrm{RSVP}}} \mathcal{F}(U'),
\]
and similarly for the Polyxan modality.

\section{Bisimulation Principles}

Two worlds are equivalent if they satisfy the same formulas under both modalities:
\[
w \sim w'
\quad\Longleftrightarrow\quad
w \Vdash A \;\Leftrightarrow\; w' \Vdash A
\quad\forall A.
\]

This bisimulation coincides with:
\begin{itemize}
\item energy–entropy equivalence on the RSVP side,
\item normal-form equivalence on the Polyxan side.
\end{itemize}

\section{Correspondence with Earlier Structures}

The geometric semantics developed here recovers the higher structures introduced throughout the monograph:

\begin{itemize}
\item curvature stratification (Ch.~13),
\item harmonic embeddings (Ch.~32),
\item confluence fibres (Ch.~25),
\item semantic galaxies (Ch.~30),
\item reconciliation normal forms (Ch.~79).
\end{itemize}

The accessibility relations coincide with the canonical flows that generate these structures.

\section{Summary}

This appendix presented:

\begin{itemize}
\item geometric frames for the two modal systems,
\item categorical interpretations of modal operators,
\item sheaf-theoretic semantics for local and global invariants,
\item bisimulation principles unifying curvature and rewriting equivalences.
\end{itemize}

This geometric foundation forms the substrate upon which the type-theoretic verification kernel is constructed in the subsequent appendix.

